\dropcap{I}n economics, marketing, and transportation research, discrete choice experiments (DCEs) are a widely used method for eliciting preferences and estimating people's willingness-to-pay (WTP) for different attributes of goods and services---travel times in transportation or air quality in environmental studies, for example. DCEs are then used to predict how people will respond to potential changes in those attributes, making them a fundamental tool for policy and business decisions \cite{mcfadden1974, train2009}.

In a typical DCE, participants choose between two or more alternatives that differ in their attributes (\autoref{fig:framework}). These choices are used to estimate the parameters of a utility function representing their underlying, stable preferences. This assumption is critical because if preferences change across the task, the estimated parameters will not reflect the true preferences of the participants, and predictions based on those estimates will be inaccurate, making DCEs unreliable for policy and business decisions. A second, fundamental assumption of DCEs is that preferences are formed by considering and comparing the whole set of attributes of each alternative. If participants do not process all of the attributes, the estimated parameters would reflect only the subset of attributes that were attended to, and predictions based on those estimates would be inaccurate.  

Both of these assumptions, however, stand in sharp contrast with decades of data from psychology and neuroscience showing that repeated choices in stable environments can lead to automatic decisions driven by habits, which make people repeat previous choices rather than evaluating deliberately the relevant attributes \cite{adamsDickinson1981, dickinson1985, dolanDayan2013, perezDickinson2020, pool2022}. Under this habit mode, changes in an otherwise stable environment can produce suboptimal decision-making \cite{hardwick2019, pool2022,henriquezJara2025}. 

As people become more familiar with a stable environment, attention progressively narrows to a subset of attributes, reducing the information considered in each choice and producing choice inertia---the tendency to repeat previous choices even when doing so is no longer optimal \cite{bouton2020, thrailkill2018}. In the animal literature, a surprising change in the environment restores attention and returns control to goal-directed behavior \cite{daw2005, lee2014, bouton2024}.

In humans, there is emerging evidence that attention plays a similar role. London Underground users often stick to familiar routes until forced by a strike to explore better alternatives \cite{larcom2017}. Participants prone to habit-like strategies focus attention on previously rewarded options while ignoring superior new alternatives---a progressive narrowing of attention, or habitual attention, whereby participants subsample information and prioritize familiar cues \cite{henriquezJara2025}. These studies show that these behavioral strategies are predicted by individual differences in psychological traits, suggesting that the tendency to rely on habitual responses reflects stable differences in how cognitive resources such as attention are deployed.

What is less clear is whether the same process unfolds in DCEs when choices are repeated over long periods, just as real-life choices are. Standard DCEs use 10--20 trials---enough to estimate preferences, but not enough to induce the overtraining that shifts behavior from deliberative to automatic in laboratory tasks. In a longer, more ecologically valid design, participants may progressively narrow attention, leading to choice inertia---a pattern that would undermine the stable-preferences assumption. This possibility is also hard to test from choices alone: a person can repeat a choice while carefully checking the attributes, or while ignoring them entirely. The observed choice is the same in both cases; only the underlying process differs. A further difficulty is establishing the causal direction between attention and preferences. Attention may reflect stable preferences---people look at what they already value---or attention may actively shape preferences---what is attended to becomes what is valued. This is an open question in the decision-making literature, and standard designs, which expose participants to a small number of novel choices, cannot distinguish the two accounts.

To test this, we used a repeated transportation-choice experiment in which participants made 150 choices among three alternatives described by cost, time, comfort, and CO\textsubscript{2} emissions, while their eye movements were recorded. This is closer to the repeated choices people face in stable everyday environments than to a short DCE. During the first 100 trials the layout was stable; at trial 101 the attribute layout changed. This design allowed us to ask whether a stable environment narrows attention, whether lower attention leads to inertia, and whether estimated WTP changes when the environment is no longer stable.

The results support this account. We identified two groups of participants: one maintained exploratory attention throughout and showed little inertia; another progressively narrowed attention, developing a more stereotyped gaze pattern and a significant increase in choice inertia during the stable phase. Critically, when a manipulation that disrupted the stability of the layout was included, attention increased for all participants---and inertia fell for those who had become automatic. This shift in attention was accompanied by an immediate change in estimated preferences. Because the spatial manipulation was exogenous---it altered the layout of the display without changing the options or their attribute values---these preference shifts are attributable to the change in attentional allocation rather than to any change in the options themselves. This identifies the causal direction: attention shapes preferences, not the other way around. Repeated choices are not repeated measurements of a fixed preference: the stability of the environment shapes the decision process itself, and disrupting that stability disrupts the preferences the model recovers. Our results raise important questions about the design of DCEs and call for a deeper integration of psychological insights into economic models of choice.